\documentclass{SibFU-docs}

% доп. пакеты
\usepackage{graphicx}
\usepackage{float}
\usepackage{hyperref}
\usepackage{caption}
\usepackage{subcaption}
\usepackage{amsmath}
% предотвартим ошибку с \Bbbk перед загрузкой пакета amssymb
\let\Bbbk\relax
\usepackage{amssymb}
\usepackage{booktabs}
\usepackage{tabularx}
\usepackage{longtable}
\usepackage{enumitem}
\usepackage{multirow}

% гиперссылки
\hypersetup{
    colorlinks=true,
    linkcolor=black,
    citecolor=blue,
    urlcolor=blue
}


% библиография
\addbibresource{report.bib}

% титульный лист
\begin{document}
\setlist[itemize]{left=1.3cm, itemsep=1pt, topsep=2pt} % для выравнивания списков
\setlist[enumerate]{left=1.3cm, itemsep=1pt, topsep=2pt}
\renewcommand{\contentsname}{}
\mycovertitle
{Гуманитарный институт}
{Кафедра прикладной информатики в искусстве и интерактивном медиа}
{Этапы жизненного цикла информационных систем}
{ст. преподаватель И.~Р.~Нигматуллин}
{Захаров И.М., Чуняев М.Д., Ситников А.В.}
 
\begin{center}
\bfseries РЕФЕРАТ
\end{center}
\noindent

Реферат по теме: «Этапы жизненного цикла информационных систем» содержит 35 страниц и 20 использованных источников.

Ключевые слова: ЖИЗНЕННЫЙ ЦИКЛ ИНФОРМАЦИОННОЙ СИСТЕМЫ, СТАДИИ ЖЦ, МОДЕЛИ ЖЦ, АНАЛИЗ ТРЕБОВАНИЙ, ПРОЕКТИРОВАНИЕ ИС, РАЗРАБОТКА И ВНЕДРЕНИЕ, ТЕСТИРОВАНИЕ И ОТЛАДКА, ЭКСПЛУАТАЦИЯ И СОПРОВОЖДЕНИЕ, СИСТЕМНОЕ АДМИНИСТРИРОВАНИЕ, ЗАВЕРШЕНИЕ ЖЦ, ОЦЕНКА ЭФФЕКТИВНОСТИ ИС

В работе проводится комплексное исследование процесса существования информационной системы от зарождения идеи до её вывода из эксплуатации.

Рассмотрены фундаментальные концепции, стандарты и модели жизненного цикла (ЖЦ) ИС, определяющие структуру и управление проектами. Детально проанализированы ключевые стадии: формирование требований и проектирование архитектуры, разработка и внедрение, а также процессы тестирования и интеграции, обеспечивающие качество системы. Особое внимание уделено самой протяжённой стадии — эксплуатации и сопровождению, раскрыты цели, задачи и практические аспекты системного администрирования. 

Завершающая часть работы посвящена методологии оценки технической, экономической и организационной эффективности ИС, которая служит основой для принятия стратегических решений о модернизации или прекращении её жизненного цикла.

Основные выводы:
\begin{enumerate}
    \item Жизненный цикл информационной системы представляет собой структурированную последовательность взаимосвязанных стадий (от анализа требований до вывода из эксплуатации), управление которыми на основе стандартов (ГОСТ 34.601-90, ISO/IEC 12207) и моделей (каскадная, гибкая) является критическим фактором успеха проекта, обеспечивая управляемость, прогнозируемость затрат и снижение рисков;
    \item Качество и успех ИС закладываются на ранних стадиях анализа требований и проектирования. Глубокое обследование предметной области, корректная спецификация функциональных и нефункциональных требований, а также выбор обоснованной архитектуры формируют технический фундамент системы, от которого напрямую зависят её надёжность, безопасность, производительность и стоимость последующего сопровождения;
    \item Стабильная и безопасная эксплуатация ИС обеспечивается комплексной деятельностью по системному администрированию, включающей управление инфраструктурой, доступом, резервным копированием, мониторингом и проактивным устранением проблем. Следование лучшим практикам (ITIL) в управлении инцидентами, проблемами и изменениями преобразует реактивную поддержку в управляемый процесс постоянного улучшения;
    \item Решение о завершении жизненного цикла ИС должно основываться на объективной, многокритериальной оценке её эффективности. Анализ технического состояния (на основе тестирования и метрик), расчёт экономических показателей (TCO, ROI) и оценка стратегической ценности для бизнеса позволяют обоснованно выбрать между модернизацией, заменой или выводом системы, превращая накопленный опыт в организационное знание для будущих проектов.
\end{enumerate}

\vspace{1em}

Рекомендации:
\begin{itemize}[label=-]
    \item При инициировании проектов по созданию ИС выбирать модель жизненного цикла (каскадную, итеративную, гибкую) осознанно, с учётом степени неопределённости требований, сроков, бюджета и зрелости команды, не допуская механического следования методологии;
    \item На стадиях анализа и проектирования уделять максимальное внимание выявлению и документированию нефункциональных требований (к безопасности, производительности, масштабируемости), так как их последующая реализация «вдогонку» является крайне затратной и часто малоэффективной;
    \item Внедрять и постоянно совершенствовать процессы системного администрирования, основанные на автоматизации рутинных операций (мониторинг, бэкапы), чётких регламентах и принципе документирования всех изменений. Это снижает операционные риски и формирует устойчивую ИТ-инфраструктуру;
    \item Планировать и проводить регулярную оценку эффективности действующих ИС, используя сбалансированный набор технических и экономических метрик. Это позволяет принимать упреждающие, а не вынужденные решения о развитии ИТ-ландшафта, оптимизируя общую стоимость владения.
\end{itemize}

\newpage

\thispagestyle{empty}
\begin{center}
\bfseries СОДЕРЖАНИЕ
\end{center}
\vspace{-6pt}
\tableofcontents
\clearpage


\section*{}
\vspace*{-2.5em}
\begin{center}\textbf{ВВЕДЕНИЕ}\end{center}
\phantomsection
\addcontentsline{toc}{section}{Введение}

Информационные системы стали ключевым элементом инфраструктуры любой современной организации, обеспечивая автоматизацию бизнес-процессов, управление данными и поддержку принятия решений. От корпоративных ERP-систем до облачных сервисов — ИС формируют основу цифровой трансформации и конкурентного преимущества. Однако создание и поддержание эффективной, надёжной и безопасной информационной системы представляет собой сложную, многогранную задачу, требующую системного подхода и грамотного управления на протяжении всего времени её существования.

В условиях быстрого технологического обновления и изменчивых бизнес-требований особенно остро встаёт проблема управления сложностью ИТ-проектов. Неструктурированный процесс разработки и внедрения зачастую приводит к срыву сроков, превышению бюджета, низкому качеству итогового продукта и, как следствие, к неудовлетворённости заказчика и пользователей. Более того, даже успешно внедрённая система требует постоянных затрат на поддержку, обновление и адаптацию, а её конечная ценность для бизнеса не всегда очевидна. В этом контексте концепция жизненного цикла информационной системы становится основополагающим инструментом, обеспечивающим дисциплину, предсказуемость и управляемость на всех этапах — от зарождения идеи до вывода системы из эксплуатации.

Современное состояние вопроса характеризуется наличием множества стандартов (ГОСТ 34.601-90, ISO/IEC 12207), методологий (водопадная, спиральная, Agile) и лучших практик (ITIL, DevOps), регламентирующих процессы ЖЦ. Однако их эффективное применение на практике требует глубокого понимания специфики каждой стадии, их взаимосвязи и вклада в общий успех проекта. Особую актуальность приобретают вопросы оценки экономической эффективности и стратегической ценности ИС, которые напрямую влияют на решения о её дальнейшем развитии или замене.

Актуальность темы обусловлена всеобщей цифровизацией, когда инвестиции в ИТ составляют значительную часть бюджета организаций. Понимание полного жизненного цикла позволяет оптимизировать эти инвестиции, минимизировать риски неудачных проектов и повысить отдачу от информационных систем. Новизна работы заключается в комплексном и последовательном рассмотрении всех этапов ЖЦ ИС — от теоретических основ и стандартов до практических аспектов администрирования и итоговой оценки эффективности, что формирует целостное представление о процессе управления ИТ-активами.

Цель работы: исследовать структуру, содержание и взаимосвязь ключевых этапов жизненного цикла информационной системы, определить их значение для успешного создания и эксплуатации ИС.

Задачи:
\begin{enumerate}
    \item Изучить понятие, модели и стандарты жизненного цикла информационной системы;
    \item Проанализировать процессы анализа требований и проектирования как фундамент будущей системы;
    \item Рассмотреть этапы разработки, тестирования, интеграции и внедрения ИС;
    \item Исследовать задачи эксплуатации, сопровождения и системного администрирования, обеспечивающие стабильную работу ИС;
    \item Проанализировать подходы к оценке эффективности ИС и методологию управляемого завершения её жизненного цикла.
\end{enumerate}

Методы: системный анализ, сравнительный анализ моделей и стандартов ЖЦ, обобщение и структурирование информации из научной литературы, нормативных документов и профессиональных практик.

\newpage
% вручную увеличить номер секции и записать в оглавление с номером,
% чтобы в теле не печатался автоматический номер, но в tableofcontents он остался
\refstepcounter{section}
\addcontentsline{toc}{section}{\protect\numberline{\thesection} Понятие жизненного цикла информационной системы и его значение}
\vspace*{-2.5em}
\begin{center}\textbf{\thesection\ ПОНЯТИЕ ЖИЗНЕННОГО ЦИКЛА ИНФОРМАЦИОННОЙ СИСТЕМЫ И ЕГО ЗНАЧЕНИЕ}\end{center}
\vspace*{-1.5em}
\ManualSubsection{Определение и стадии жизненного цикла}

Жизненный цикл информационной системы (ЖЦ ИС) представляет собой период времени от момента принятия решения о создании системы до её полного вывода из эксплуатации. Это концепция, позволяющая управлять сложным процессом создания и использования ИС как последовательностью взаимосвязанных этапов, каждый из которых имеет свои цели, результаты, методы и ответственных лиц.

Ключевые стадии ЖЦ образуют универсальную последовательность:

\begin{enumerate}
\item Формирование требований: анализ предметной области, планирование, построение моделей «As-Is» и «To-Be»;
\item Проектирование: определение архитектуры, функций, интерфейсов и требований к компонентам;
\item Реализация: непосредственное программирование и создание компонентов;
\item Внедрение: установка, настройка, обучение пользователей и промышленный запуск;
\item Эксплуатация и сопровождение: повседневная работа, поддержка, исправление ошибок, адаптация;
\item Вывод из эксплуатации: миграция данных, остановка и утилизация системы.
\end{enumerate}

\vspace{1em}

Формализация ЖЦ позволяет превратить сложную задачу создания ИС в управляемый процесс с чёткими контрольными точками, что критически важно для планирования сроков, бюджета и ресурсов. \cite{gost3460190, ISO12207_2017, Dushin2006}

\ManualSubsection{Классификация процессов жизненного цикла по ISO/IEC 12207}

Стандарт ISO/IEC 12207 рассматривает жизненный цикл не только как набор стадий, но и как совокупность процессов, которые делятся на три основные группы:
\begin{itemize}[label=-]
\item Основные процессы: непосредственно связаны с созданием и использованием системы (заказ, поставка, разработка, эксплуатация, сопровождение);
\item Вспомогательные процессы: обеспечивают качественное выполнение основных (документирование, управление конфигурацией, обеспечение качества, верификация, аттестация, совместная оценка, аудит);
\item Организационные процессы: направлены на управление проектом и улучшение самих процессов ЖЦ (управление, создание инфраструктуры, усовершенствование, обучение).
\end{itemize}

\vspace{1em}

Такой процессно-ориентированный подход обеспечивает гибкость и позволяет адаптировать ЖЦ под конкретные нужды проекта, фокусируясь на качестве и управлении. \cite{ISO12207_2017}

\ManualSubsubsection{Взаимосвязь стадий и процессов}

Стадии (этапы) задают временную шкалу и ключевые результаты проекта, в то время как процессы описывают, какие действия и каким образом должны выполняться для достижения этих результатов. Например, процесс «управление качеством» или «верификация» активен на протяжении нескольких стадий, обеспечивая сквозной контроль.

\ManualSubsection{Эволюция моделей жизненного цикла: от каскадных к гибким}

Исторически первой и наиболее формализованной является каскадная (водопадная) модель, регламентированная в ГОСТ 34.601-90. Её ключевой принцип — строго последовательное выполнение стадий, где переход к следующему этапу возможен только после полного завершения и документального оформления предыдущего. Эта модель обеспечивает высокую предсказуемость и контроль, но критически уязвима к изменениям требований, которые могут возникнуть на поздних стадиях.

В ответ на эту проблему появились итеративные и инкрементальные модели, такие как Rational Unified Process (RUP) и Microsoft Solution Framework (MSF). Они разбивают разработку на циклы (итерации), каждый из которых проходит мини-жизненный цикл и завершается выпуском рабочего инкремента системы. Это позволяет постепенно уточнять требования и раньше получать обратную связь от заказчика.

Дальнейшим развитием стали гибкие методологии (Agile), такие как Scrum и Extreme Programming (XP). Их философия отражена в Agile-манифесте: люди и взаимодействие важнее процессов и инструментов, работающий продукт важнее исчерпывающей документации, сотрудничество с заказчиком важнее согласования условий контракта, готовность к изменениям важнее следования первоначальному плану. В современных условиях проекты всё чаще используют гибридные подходы, сочетающие структурированность каскадного планирования с гибкостью Agile-итераций. \cite{gost3460190, sommerville2016se, schwaber2020scrum, pressman2014se}

\ManualSubsubsection{Спиральная модель как управление рисками}

Отдельного внимания заслуживает спиральная модель Боэма, которая делает явный акцент на анализе и минимизации рисков на каждой итерации. Каждый виток спирали включает этапы: определение целей и ограничений, оценку и уменьшение рисков, разработку и тестирование, а также планирование следующей итерации. Эта модель эффективна для крупных, инновационных и рискованных проектов.

\ManualSubsection{Системное значение жизненного цикла для организации}

Понимание и грамотное применение принципов ЖЦ — это не техническая формальность, а основа управленческой зрелости в сфере ИТ. Оно позволяет:

\begin{enumerate}

\item Снижать риски и стоимость ошибок. Ошибка, обнаруженная на стадии формирования требований, исправляется в десятки раз дешевле, чем та же ошибка, найденная на стадии внедрения или эксплуатации;
\item Обеспечивать долгосрочную экономическую эффективность. Акцент на стадиях эксплуатации и сопровождения (которые поглощают 60-80 процентов общих затрат) заставляет планировать систему с учётом стоимости владения, а не только стоимости разработки;
\item Накопить организационные знания. Формализация процессов и документирование опыта завершённых проектов создают базу для повышения эффективности будущих инициатив;
\item Улучшать коммуникацию. Чёткое распределение ролей, ответственности и результатов между заказчиками, аналитиками, разработчиками и пользователями минимизирует недопонимание.
\end{enumerate}

\vspace{1em}

Таким образом, жизненный цикл информационной системы является фундаментальной концепцией, связывающей технические, управленческие и экономические аспекты создания и использования ИТ-активов в организации. \cite{gost3460190, ISO12207_2017, OGC_ITIL_2007}


\newpage
\refstepcounter{section}
\addcontentsline{toc}{section}{\protect\numberline{\thesection}Анализ требований и проектирование структуры будущей системы}
\begin{center}\textbf{\thesection\ АНАЛИЗ ТРЕБОВАНИЙ И ПРОЕКТИРОВАНИЕ СТРУКТУРЫ БУДУЩЕЙ СИСТЕМЫ}\end{center}
\vspace*{-1.5em}
\ManualSubsection{Глубокий анализ как основа успешного проекта}

Этап анализа требований — это процесс перевода расплывчатых пожеланий заказчика в чёткую, полную, непротиворечивую и проверяемую спецификацию. Его основная задача — дать однозначный ответ на вопрос: «Что должна делать система?». Работа начинается с погружения в предметную область. Аналитики строят модель текущего состояния бизнес-процессов («As-Is»), что позволяет выявить существующие узкие места, неэффективности и рутинные операции, подлежащие автоматизации. На основе этого анализа и стратегических целей организации формируется целевая модель («To-Be») — видение оптимизированных процессов, поддерживаемых новой ИС.

\ManualSubsubsection{Методы обследования предметной области и выявления требований}

Для сбора информации используются различные методы, выбор которых зависит от контекста проекта:

\begin{itemize}[label=-]
\item Интервью (индивидуальные и групповые) позволяют получить глубокие качественные данные от ключевых пользователей и экспертов;
\item Анкетирование эффективно для массового сбора мнений и количественных данных;
\item Наблюдение (пассивное или активное) за работой сотрудников помогает выявить неочевидные рутинные операции и реальные сценарии работы;
\item Анализ документов (регламентов, отчётов, форм) и существующих информационных систем;
\item Семинары по совместному проектированию (JAD-сессии), где все заинтересованные стороны (бизнес-пользователи, аналитики, разработчики) совместно прорабатывают требования и архитектурные решения, что значительно ускоряет процесс и повышает согласованность.
\end{itemize}

\vspace{1em}

\ManualSubsubsection{Классификация и приоритезация требований}

Выявленные требования структурируются и классифицируются:

\begin{itemize}[label=-]
\item Бизнес-требования: высокоуровневые цели организации (например, «сократить цикл обработки заказа на 30 процентов»);
\item Требования пользователей (User Needs): задачи, которые пользователи должны решать с помощью системы, часто описываемые в виде сценариев или пользовательских историй (User Stories);
\item Функциональные требования: детальное описание поведения системы, её функций и реакций на входные данные. Формализуются с помощью вариантов использования (Use Cases), которые описывают взаимодействие актора (пользователя или внешней системы) с ИС для достижения конкретной цели;
\item Нефункциональные требования (или требования к качеству): определяют, как система должна выполнять свои функции.
\end{itemize}

\vspace{1em}

Итогом этапа является Спецификация требований к программному обеспечению (SRS) — документ, который служит юридическим и техническим контрактом между заказчиком и разработчиком, а также основой для проектирования и тестирования. \cite{kotonya1998requirements, ieee830, ISO25010_2011}

\ManualSubsection{Архитектурное и детальное проектирование системы}

Этап проектирования трансформирует «что» в «как». На его основе создаётся целостная техническая модель системы, которая определяет, как будут реализованы все требования.

\ManualSubsubsection{Проектирование архитектуры системы}
Это высокоуровневый выбор структуры и технологий. Определяется архитектурный стиль:

\begin{itemize}[label=-]
\item Клиент-сервер (2-уровневая): разделение на клиентское приложение и сервер БД;
\item Многоуровневая (N- tier) архитектура: разделение на уровень представления (UI), уровень бизнес-логики (сервер приложений) и уровень данных (БД). Повышает масштабируемость и безопасность;
\item Сервис-ориентированная архитектура (SOA) и микросервисы: система строится как набор слабо связанных, независимо развёртываемых сервисов, общающихся по сети. Обеспечивает высокую гибкость и скорость разработки;
\item Веб-архитектура: клиент (браузер) взаимодействует с веб-сервером по протоколам HTTP/HTTPS.
\end{itemize}

\vspace{1em}

Архитектурное проектирование также включает выбор технологического стека (языки программирования, фреймворки, СУБД), проектирование интеграционных интерфейсов (API) и определение стратегии развёртывания.

\ManualSubsubsection{Проектирование данных}

Создаётся детальная модель данных, которая проходит этапы:

\begin{enumerate}
    \item Концептуальное проектирование: создание ER-диаграмм, отображающих сущности предметной области и связи между ними;
    \item Логическое проектирование: отображение концептуальной модели в структуры конкретной СУБД (таблицы, столбцы, первичные и внешние ключи) с применением нормализации для минимизации избыточности;
    \item Физическое проектирование: оптимизация для производительности (индексы, партиционирование, денормализация), определение стратегий резервного копирования и восстановления.
\end{enumerate}

\ManualSubsubsection{Проектирование пользовательского интерфейса (UI/UX)}

Разрабатываются макеты (wireframes) и прототипы экранов системы. Цель — создать интуитивно понятный, эффективный и эстетичный интерфейс, который минимизирует ошибки пользователя и время на обучение. Применяются принципы юзабилити, проводится тестирование прототипов на фокус-группах.

\ManualSubsubsection{Спецификация системных требований}

На основе нефункциональных требований формируются чёткие системные требования, которые определяют необходимую инфраструктуру:

\begin{itemize}[label=-]
    \item Аппаратные требования к серверам (CPU, RAM, диск);
    \item Требования к сетевой инфраструктуре (пропускная способность, топология).
    \item Требования к системному ПО (операционные системы, веб-серверы, middleware);
    \item Требования к безопасности (межсетевые экраны, системы обнаружения вторжений).
\end{itemize}

Итогом проектирования является Технический проект (ТП) — комплект документации, содержащий все архитектурные схемы, модели данных, спецификации интерфейсов и системные требования, достаточные для передачи в разработку. \cite{sommerville2016se, pressman2014se, gost3460190}

\newpage
\refstepcounter{section}
\addcontentsline{toc}{section}{\protect\numberline{\thesection}Этапы разработки и внедрения информационной системы}
\begin{center}\textbf{\thesection\ ЭТАПЫ РАЗРАБОТКИ И ВНЕДРЕНИЯ ИНФОРМАЦИОННОЙ СИСТЕМЫ}\end{center}
\vspace*{-1.5em}
\ManualSubsection{Реализация (разработка) программных компонентов}

На стадии реализации технический проект превращается в рабочий код. Современная разработка — это не просто написание программ, а управляемый инженерный процесс.

\ManualSubsubsection{Организация процесса разработки}

Ключевые практики:

\begin{itemize}[label=-]
    \item Версионный контроль (Git): обязателен для совместной работы, отслеживания изменений и возможности отката;
    \item Непрерывная интеграция (CI): автоматическая сборка и базовое тестирование кода при каждом изменении в репозитории. Позволяет быстро обнаруживать проблемы интеграции;
    \item Парное программирование и ревью кода: коллегиальная проверка кода для повышения его качества, выявления ошибок и распространения знаний в команде;
    \item Следование стандартам кодирования и принципам проектирования (SOLID, DRY, KISS): обеспечивает создание сопровождаемого, модульного и гибкого кода.
\end{itemize}

Разработка часто ведётся по модульному принципу, когда разные команды или разработчики параллельно работают над отдельными компонентами системы. \cite{sommerville2016se, pressman2014se}

\ManualSubsubsection{Интеграция компонентов}

По мере готовности модули начинают объединяться. Для управления зависимостями и сборкой используются инструменты вроде Maven, Gradle, npm. Проводится интеграционное тестирование, которое проверяет корректность взаимодействия между модулями, правильность передачи данных через API и общую работоспособность подсистем. Стратегии интеграции (нисходящая, восходящая, «Большой взрыв») выбираются исходя из архитектуры проекта.

\ManualSubsection{Всестороннее тестирование системы}

Тестирование — это не этап, а сквозная деятельность, пик которой приходится на период после сборки системы. Его цель — верификация соответствия системы спецификации требований (SRS) и валидация (удовлетворяет ли она потребностям пользователя).

\ManualSubsubsection{Уровни и виды тестирования}

\begin{itemize}[label=-]
    \item Модульное (Unit) тестирование: проверка наименьших неделимых единиц кода (функций, методов) изолированно. Часто выполняется самими разработчиками с использованием фреймворков (JUnit, pytest);
    \item Интеграционное тестирование: проверка взаимодействия модулей, интеграции с внешними системами (платёжные шлюзы, SMS-сервисы);
    \item Системное тестирование: проверка системы в целом;
    \item Приёмочное тестирование (User Acceptance Testing — UAT): финальная проверка, которую проводят представители заказчика или реальные пользователи в среде, максимально приближенной к промышленной. Успешное UAT — формальный знак готовности системы к внедрению. \cite{Kulikov2017Testing, ISO29119-1-2022}
\end{itemize}

\ManualSubsection{Планирование и осуществление внедрения}

Внедрение — критический переход от разработки к реальной эксплуатации. Это организационно-технический процесс, требующий тщательного планирования.

\ManualSubsubsection{Подготовительные мероприятия}

\begin{itemize}[label=-]
    \item Подготовка инфраструктуры: развёртывание и настройка серверов, сетевого оборудования, систем хранения данных (СХД), установка необходимого системного ПО;
    \item Миграция данных: один из самых сложных и рискованных процессов. Включает извлечение данных из старых систем (legacy), их очистку (data cleansing), преобразование (transformation) и загрузку (loading — ETL) в новую ИС. Требует разработки специальных скриптов и проведения нескольких репетиций;
    \item Обучение пользователей и администраторов: разработка учебных материалов, проведение тренингов и инструктажей. Квалифицированные пользователи — ключ к успешному внедрению.
\end{itemize}

\ManualSubsubsection{Стратегии внедрения и переходный период}

Выбор стратегии зависит от критичности системы и допустимого риска:

\begin{itemize}[label=-]
    \item Прямой переход («Big Bang»): одномоментный отказ от старой системы и запуск новой. Рискованный, но быстрый;
    \item Параллельный запуск: новая и старая системы работают одновременно какое-то время. Снижает риск, но удваивает нагрузку на персонал;
    \item Пилотное (поэтапное) внедрение: система запускается сначала в одном подразделении или для одного процесса. После успешной обкатки и устранения проблем масштабируется на всю организацию. Наиболее безопасная и рекомендуемая стратегия.
\end{itemize}

После запуска начинается период гиперподдержки, когда разработчики и внедренцы находятся в максимальной готовности для оперативного решения возникающих у пользователей проблем. \cite{pressman2014se}



\newpage
\refstepcounter{section}
\addcontentsline{toc}{section}{\protect\numberline{\thesection}Тестирование, отладка и интеграция функциональных компонентов}
\begin{center}\textbf{\thesection\ ТЕСТИРОВАНИЕ, ОТКЛАДКА И ИНТЕГРАЦИЯ ФУНКЦИОНАЛЬНЫХ КОМПОНЕНТОВ}\end{center}
\vspace*{-1.5em}

\ManualSubsection{Системный подход к обеспечению качества через тестирование}

Тестирование в современной разработке — это не поиск ошибок по окончании кодирования, а процесс, интегрированный в жизненный цикл и направленный на предотвращение дефектов. Стандарт ISO/IEC/IEEE 29119 формализует процессы, документацию и техники тестирования, подчёркивая его систематический характер.

\ManualSubsubsection{Статическое и динамическое тестирование}

\begin{itemize}[label=-]
    \item Статическое тестирование: анализ артефактов без запуска программы. Включает инспекции и рецензирование требований, архитектурных схем, кода (code review). Позволяет выявить до 60 процентов дефектов на ранних и дешёвых стадиях;
    \item Динамическое тестирование: проверка поведения программы в ходе её выполнения. Все виды тестирования, связанные с запуском кода (модульное, интеграционное, системное), являются динамическими.
\end{itemize}

\ManualSubsubsection{Автоматизация тестирования}

Краеугольный камень DevOps-культуры. Автоматизированные тесты (unit, интеграционные, end-to-end) обеспечивают:

\begin{itemize}[label=-]
    \item Быструю обратную связь для разработчиков;
    \item Возможность регрессионного тестирования — проверки, что новые изменения не сломали существующую функциональность;
    \item Непрерывное обеспечение качества в процессе частых выпусков.
\end{itemize}

Написание тестовых сценариев рассматривается как такая же важная часть разработки, как и написание основного кода. \cite{Kulikov2017Testing, ISO29119-1-2022, Murphy2016SRE}

\ManualSubsection{Отладка как процесс исследования системы}

Отладка (debugging) — это деятельность по локализации, анализу и устранению причин дефектов, обнаруженных в ходе тестирования или эксплуатации.

\ManualSubsubsection{Инструменты и методы отладки}

\begin{itemize}[label=-]
    \item Интерактивные отладчики (debuggers): интегрированы в среды разработки (IDE). Позволяют выполнять код по шагам, устанавливать точки останова, инспектировать значения переменных в режиме реального времени, анализировать стек вызовов. Незаменимы для анализа логических ошибок;
    \item Ведение журналов (логирование): стратегическое размещение операторов записи в журнал (logging) в критичных точках программы. Анализ логов — основной метод диагностики проблем в промышленной среде, особенно для распределённых систем, где воспроизвести ошибку в отладчике сложно.
    \item Профилирование (profiling): измерение потребления ресурсов (время выполнения, использование памяти) для поиска «узких мест» в производительности.
\end{itemize}

\ManualSubsubsection{Анализ инцидентов в промышленной среде}

В эксплуатации отладка часто связана с расследованием инцидентов. Используется централизованный сбор логов (например, в ELK-стек или Splunk), система мониторинга (Zabbix, Prometheus) и трейсинг распределённых транзакций (Jaeger, Zipkin). Это позволяет восстановить цепочку событий, приведшую к сбою, даже если она затрагивает десятки микросервисов. \cite{Murphy2016SRE}

\ManualSubsection{Интеграция компонентов в единое целое}

Интеграция — это процесс сборки отдельных программных модулей в согласованно работающую систему.

\ManualSubsubsection{Стратегии интеграции}
\begin{enumerate}
    \item Нисходящая интеграция (Top-Down): сборка начинается с модулей верхнего уровня (например, пользовательского интерфейса), а для отсутствующих нижних модулей используются временные «заглушки» (stubs);
    \item Восходящая интеграция (Bottom-Up): сначала интегрируются нижние, фундаментальные модули (например, доступ к данным), а для верхних уровней создаются «драйверы» (drivers).
    \item Сэндвич-стратегия (гибридная): комбинация нисходящего и восходящего подходов;
    \item Большой взрыв (Big Bang): все модули разрабатываются независимо и интегрируются одновременно в самом конце. Крайне рискованна, так как приводит к «интеграционному хаосу».
\end{enumerate}

\ManualSubsubsection{Непрерывная интеграция как культура}

Практика Continuous Integration (CI) — это реализация частой (несколько раз в день) автоматической интеграции изменений от всех разработчиков в общую основную ветку. Каждое слияние сопровождается автоматизированной сборкой и прогоном набора тестов. Это позволяет практически мгновенно обнаруживать интеграционные проблемы, существенно снижая стоимость их исправления. CI является фундаментальной практикой для Agile- и DevOps-команд. \cite{sommerville2016se}

\ManualSubsubsection{Интеграционное тестирование}

Проверяет корректность взаимодействия между модулями, подсистемами и внешними системами. Фокусируется на интерфейсах, протоколах обмена данными, обработке ошибок в межмодульном взаимодействии. Часто требует создания сложного тестового окружения, имитирующего работу смежных систем.


\newpage
\refstepcounter{section}
\addcontentsline{toc}{section}{\protect\numberline{\thesection}Эксплуатация, сопровождение и модернизация системы}
\begin{center}\textbf{\thesection\ ЭКСПЛУАТАЦИЯ, СОПРОВОЖДЕНИЕ И МОДЕРНИЗАЦИЯ СИСТЕМЫ}\end{center}
\vspace*{-1.5em}
\ManualSubsection{Системное администрирование как основа стабильной эксплуатации ИС}

После успешного внедрения информационная система переходит в самую длительную и ресурсоёмкую фазу своего существования — эксплуатацию. На этом этапе основная цель смещается с создания системы на обеспечение её стабильной, безопасной и эффективной работы для достижения бизнес-целей организации. Ключевую роль в этом играет системное администрирование — комплекс практических действий, направленных на поддержание жизнеспособности ИС в реальных условиях. Администратор действует на стыке всех уровней архитектуры ИС: от аппаратного (серверы, сети, системы хранения) и системного (ОС, гипервизоры, СУБД) до прикладного (бизнес-приложения) и пользовательского, обеспечивая их слаженное взаимодействие. \cite{Murphy2016SRE, studfile_admin_ib_2025}

Главными целями системного администрирования являются: обеспечение доступности систем в нужное для бизнеса время; сохранение целостности и надёжности данных; поддержание необходимого уровня производительности; соблюдение требований информационной безопасности; а также обеспечение управляемости и масштабируемости ИС для её будущего развития.

\ManualSubsection{Ключевые направления деятельности системного администратора}

Деятельность системного администратора структурирована вокруг нескольких основных блоков задач, каждый из которых критически важен для устойчивой работы ИС.

\ManualSubsubsection{Управление инфраструктурой и ресурсами}

Это фундаментальная задача, включающая развёртывание, настройку и обслуживание всего аппаратно-программного комплекса: физических и виртуальных серверов, сетевого оборудования (коммутаторов, маршрутизаторов, межсетевых экранов), систем хранения данных (СХД), а также инженерной инфраструктуры (электропитание, охлаждение, противопожарные системы). Администратор планирует ресурсы (вычислительные мощности, оперативную память, дисковое пространство, сетевую пропускную способность) с запасом на рост нагрузки и осуществляет своевременную модернизацию или замену устаревающего оборудования. Без корректно настроенной и стабильной инфраструктуры даже самое совершенное прикладное программное обеспечение не сможет функционировать эффективно.

\ManualSubsubsection{Управление доступом и обеспечение информационной безопасности}

Администратор управляет жизненным циклом учётных записей пользователей и служб: создаёт их, назначает права в соответствии с принципом минимально необходимых привилегий, блокирует и удаляет при увольнении сотрудников. Эта деятельность тесно переплетается с задачами информационной безопасности. Помимо управления правами, администратор обеспечивает своевременное применение обновлений безопасности (патчей) для операционных систем и приложений, проводит сегментацию сети для ограничения доступа к критическим ресурсам, настраивает и анализирует журналы аудита (логи) для выявления подозрительных активностей и участвует в расследовании инцидентов безопасности. Защита ИС строится не только на формальных политиках, но и на ежедневных рутинных действиях администратора, где ошибки могут привести к серьёзным утечкам данных или простоям. \cite{studfile_admin_ib_2025}

\ManualSubsubsection{Организация резервного копирования и восстановления}

Резервное копирование — основной инструмент защиты от потери данных в результате сбоев оборудования, человеческих ошибок, вредоносных атак (например, программ-вымогателей) или катастроф. Задача администратора — разработать и реализовать надёжную стратегию резервного копирования. Это включает определение критичных данных и частоты их копирования, выбор типов бэкапов (полные, инкрементальные, дифференциальные), следование правилу 3-2-1 (три копии данных, на двух разных типах носителей, одна копия географически удалённо), а также обеспечение защиты самих резервных копий от повреждения и несанкционированного доступа. Крайне важно регулярно тестировать процедуру восстановления, так как непригодная резервная копия хуже, чем её отсутствие. \cite{veeam_backup_recovery_2024}

\ManualSubsubsection{Мониторинг, диагностика и управление инцидентами}

Постоянный мониторинг состояния всех компонентов ИС (доступность, загрузка процессора и памяти, свободное место на дисках, сетевой трафик) позволяет выявлять проблемы до того, как они повлияют на пользователей. Для этого используются специализированные системы мониторинга, которые собирают метрики, логи и генерируют оповещения (алерты).

В рамках лучших практик ITIL администратор участвует в процессах управления инцидентами и проблемами. Инцидент — это любое отклонение от штатной работы сервиса (например, недоступность системы). Цель — как можно быстрее восстановить обслуживание, даже если временными мерами. Проблема — это коренная причина одного или нескольких инцидентов (например, ошибка в коде или неисправный модуль памяти). Цель управления проблемами — найти и устранить эту первопричину, чтобы инциденты не повторялись. Таким образом, администратор не только «тушит пожары», но и системно улучшает устойчивость ИС. \cite{OGC_ITIL_2007}

\ManualSubsubsection{Профилактическое обслуживание и документирование}
Проактивная профилактика — залог предсказуемой работы. К ней относятся плановые регламентные работы: перезагрузка серверов для применения обновлений, проверка целостности данных, очистка временных файлов, тестирование систем резервного копирования и аварийного восстановления, контроль параметров среды в серверной (температура, влажность). Отдельная критическая задача — документирование: ведение актуальной документации по архитектуре, конфигурациям, сетевым схемам, процедурам восстановления. Если знания остаются только в голове администратора, его уход или ошибка становятся серьёзным риском для бизнеса.

\ManualSubsection{Сопровождение и эволюционная модернизация ИС}

Сопровождение программного обеспечения — это процесс модификации системы после сдачи её в эксплуатацию для исправления ошибок, улучшения характеристик или адаптации к изменяющемуся окружению. Согласно стандарту ISO/IEC 14764, выделяют три основных типа работ по сопровождению:

\begin{enumerate}
    \item Корректирующее: устранение выявленных в эксплуатации дефектов и ошибок;
    \item Адаптивное: доработка системы для обеспечения её работы в новой или изменённой среде (например, при обновлении версии операционной системы или СУБД);
    \item Совершенствующее: улучшение функциональных возможностей или производительности системы в ответ на новые или изменившиеся требования пользователей. Именно этот тип работ составляет основную часть активности после первоначального внедрения и часто реализуется в виде выпусков новых версий или обновлений. \cite{ISO12207_2017, pressman2014se}
\end{enumerate}

Когда накопленные изменения становятся слишком велики, затраты на сопровождение резко возрастают, а технологический стек устаревает, встаёт вопрос о модернизации (реинжиниринге). Это может быть как масштабный рефакторинг кода и перенос на новую платформу, так и полная замена системы. Администратор, обладающий глубокими знаниями о текущей архитектуре и её «болевых точках», играет ключевую роль в планировании такой модернизации, оценке рисков и обеспечении непрерывности бизнес-процессов в переходный период.

\newpage
\refstepcounter{section}
\addcontentsline{toc}{section}{\protect\numberline{\thesection}Завершение жизненного цикла и оценка эффективности информационной системы}
\begin{center}\textbf{\thesection\ ЗАВЕРШЕНИЕ ЖИЗНЕННОГО ЦИКЛА И ОЦЕНКА ЭФФЕКТИВНОСТИ ИНФОРМАЦИОННОЙ СИСТЕМЫ}\end{center}
\vspace*{-1.5em}

\ManualSubsection{Завершающая стадия жизненного цикла: сущность и необходимость}

Завершение жизненного цикла — закономерный и неизбежный этап существования любой информационной системы. Он наступает, когда дальнейшая эксплуатация ИС становится экономически нецелесообразной или технически невозможной. Причины могут быть различными: моральное и физическое устаревание («устаревание»), кардинальное изменение бизнес-процессов, неподъёмные затраты на поддержку устаревших технологий, появление на рынке принципиально новых, более эффективных решений или смена стратегических приоритетов организации.

Эта стадия выходит далеко за рамки простого «выключения» сервера. Это планомерный, управляемый процесс, который включает оценку накопленного опыта, архивирование данных, перевод пользователей на новые решения (миграцию) и формальное прекращение эксплуатации. Решение о завершении жизненного цикла должно основываться не на интуиции, а на системной аналитике, использующей данные эксплуатации, результаты тестирования и экономические расчёты. \cite{OGC_ITIL_2007, gost3460190}

\ManualSubsection{Роль тестирования и отладки в оценке технического состояния ИС}

На завершающем этапе процессы тестирования, отладки и интеграции меняют свою направленность. Их цель смещается с развития системы на формальную верификацию её текущего состояния и сбор объективных данных для принятия управленческого решения.

\ManualSubsubsection{Тестирование как инструмент аудита}

Проводятся специальные оценочные испытания, часто повторяющие тесты, выполнявшиеся при внедрении. Регрессионное тестирование проверяет, сохранила ли система ключевую функциональность. Нагрузочное и стресс-тестирование выявляют реальные пределы производительности и запас прочности. Приёмочное тестирование (UAT), проводимое с ключевыми пользователями, фиксирует фактические показатели удобства, времени выполнения операций и количества ошибок в реальных сценариях. Эти результаты дают количественную оценку технического износа системы. \cite{ISO29119-1-2022, Kulikov2017Testing}

\ManualSubsubsection{Отладка и мониторинг для анализа надёжности}

В период зрелой эксплуатации отладка фокусируется на анализе накопленных данных. Изучается статистика инцидентов, среднее время наработки на отказ (MTBF) и среднее время восстановления (MTTR). Анализ логов (журналов) работы системы, особенно с применением инструментов для корреляции событий и обнаружения аномалий, позволяет выявить скрытые, системные проблемы, которые накапливались годами. Эта информация критически важна для оценки сопровождаемости и стоимости будущей поддержки. \cite{Murphy2016SRE}

\ManualSubsubsection{Анализ интеграционных связей}

Перед выводом системы из эксплуатации необходим тщательный аудит всех её интеграционных интерфейсов (API, обмен файлами, доступ к общим базам данных). Требуется составить полную карту зависимостей: какие внешние и внутренние системы потребляют данные или сервисы от данной ИС. Это определяет сложность и риски предстоящей миграции. Проводится интеграционное тестирование сценариев отключения, чтобы убедиться, что остановка одной системы не вызовет каскадных сбоев в других.

\ManualSubsection{Комплексная оценка эффективности информационной системы}

Оценка эффективности — это стержень, вокруг которого строится обоснование для завершения жизненного цикла. Она должна быть многомерной, сопоставляя первоначальные цели проекта с фактическими результатами за весь период эксплуатации.

\ManualSubsubsection{Оценка технической и эксплуатационной эффективности}

Оценивается соответствие системы требованиям к качеству. Используется модель, описанная в стандарте ISO/IEC 25010, которая включает такие характеристики, как функциональная полнота, производительность, надёжность, безопасность и сопровождаемость. Фактические показатели (время отклика, процент доступности, количество критических инцидентов) сравниваются с плановыми значениями, зафиксированными в техническом задании. Система, постоянно не соответствующая требованиям по ключевым параметрам, является кандидатом на замену. \cite{ISO25010_2011}

\ManualSubsubsection{Оценка экономической эффективности}

Это наиболее весомый для бизнеса аспект. Проводится анализ совокупной стоимости владения (TCO, Total Cost of Ownership), которая включает все прямые и косвенные затраты за весь жизненный цикл: капитальные затраты на разработку/покупку и внедрение, а также операционные затраты на сопровождение, обновления, администрирование, электроэнергию и т.д.

Полученные значения сопоставляются с достигнутыми выгодами (Benefits). Ключевой метрикой является возврат на инвестиции (ROI — Return on Investment), которая рассчитывает, какую прибыль или экономию принесла система по отношению к вложенным в неё средствам. Если рост затрат на сопровождение устаревающей системы начинает превышать получаемую от неё пользу или ожидаемую выгоду от модернизации, экономическое обоснование для её дальнейшего существования исчезает. \cite{Botchkarev2011, Anisiforov2014}

\ManualSubsubsection{Оценка стратегической и организационной эффективности}

Оценивается, насколько ИС поддерживает текущие и, что важно, стратегические цели бизнеса. Способствует ли она повышению конкурентоспособности, выходу на новые рынки, улучшению качества обслуживания клиентов? Также анализируется организационное влияние: повысила ли система прозрачность процессов, скорость принятия решений, удовлетворённость сотрудников? Система, ставшая тормозом для организационного развития, подлежит замене, даже если технически ещё работоспособна.

\ManualSubsection{Принятие решения и организация вывода системы из эксплуатации}

На основе комплексной оценки готовится итоговое заключение об эффективности ИС, в котором формулируются возможные сценарии:

\begin{enumerate}
    \item Продление эксплуатации с минимальными доработками (если система удовлетворительно эффективна и затраты на замену высоки);
    \item Модернизация (если ядро системы жизнеспособно, но требуется значительное обновление компонентов);
    \item Поэтапная замена на новое решение;
    \item Полный вывод из эксплуатации (если функция утратила актуальность или система экономически нецелесообразна).
\end{enumerate}

Выбранный сценарий детализируется в план вывода из эксплуатации. Этот план включает: график отключения функций, процедуры окончательного архивирования и консервации данных (в соответствии с законодательными и регуляторными требованиями), план миграции пользователей на новые системы, мероприятия по информированию всех заинтересованных сторон и минимизации рисков для бизнеса. Важным итогом является документирование извлечённых уроков — успешных практик и допущенных ошибок, которые становятся ценным организационным знанием для будущих проектов. \cite{OGC_ITIL_2007, gost3460190}

Таким образом, завершение жизненного цикла — это не конец, а важная точка в процессе непрерывного развития корпоративной информационной среды, позволяющая высвободить ресурсы для более перспективных и эффективных решений.

\newpage
\section*{}
\vspace*{-2.5em}
\begin{center}\textbf{ЗАКЛЮЧЕНИЕ}\end{center}
\phantomsection
\addcontentsline{toc}{section}{Заключение}

Жизненный цикл информационной системы представляет собой непрерывный и управляемый процесс, охватывающий все этапы её существования. Исследование показало, что эффективность ИС в конечном итоге определяется не только качеством программного кода, но и тем, насколько системно и профессионально организована работа на каждой стадии этого цикла — от первоначального замысла до окончательного вывода из эксплуатации.

Фундаментом успеха любого ИТ-проекта являются ранние стадии — анализ требований и проектирование. Именно здесь закладываются архитектурные решения, определяющие такие ключевые характеристики системы, как масштабируемость, безопасность, производительность и сопровождаемость. Ошибки и недоработки, допущенные на этом этапе, обходятся исключительно дорого на последующих стадиях, подчёркивая критическую важность глубокого погружения в предметную область и тщательной проработки технической спецификации.

Стадии разработки, тестирования и внедрения переводят проект из области проектирования в область практической реализации. Современные подходы, такие как непрерывная интеграция, автоматизированное тестирование и гибкие стратегии внедрения, позволяют минимизировать риски, повысить качество продукта и обеспечить его соответствие ожиданиям пользователей. Однако передача системы в эксплуатацию — это не конечная точка, а начало нового, самого длительного этапа её жизни.

Эксплуатация и сопровождение ИС, обеспечиваемые комплексной деятельностью по системному администрированию, являются залогом её долгосрочной стабильности и ценности для бизнеса. Управление инфраструктурой, доступом, мониторинг, резервное копирование и проактивное устранение проблем формируют среду, в которой система может надёжно выполнять свои функции. Принципы ITIL, разделяющие реактивное восстановление сервиса (управление инцидентами) и поиск коренных причин сбоев (управление проблемами), превращают поддержку из затратной функции в процесс постоянного улучшения.

Завершающая стадия жизненного цикла, основанная на всесторонней оценке эффективности, знаменует переход от эксплуатации к стратегическому управлению ИТ-портфелем. Объективный анализ технического состояния, экономических показателей (TCO, ROI) и стратегического соответствия позволяет принять обоснованное решение о модернизации, замене или выводе системы. Этот процесс не только завершает жизнь конкретной ИС, но и генерирует бесценные организационные знания, снижая риски и повышая эффективность будущих проектов.

Таким образом, этапы жизненного цикла информационной системы — это не изолированные друг от друга фазы, а элементы единого, целостного процесса управления ИТ-активом. Только комплексный подход, учитывающий методологические, технические, экономические и организационные аспекты на каждом этапе, позволяет создавать, поддерживать и развивать информационные системы, которые действительно становятся драйверами роста и эффективности организации в цифровую эпоху.

\newpage
\begin{center}\textbf{СПИСОК СОКРАЩЕНИЙ}\end{center}
\phantomsection
\addcontentsline{toc}{section}{Список сокращений}

\vspace{1em}

{\small
\noindent
\begin{tabular}{|p{0.15\textwidth}|p{0.80\textwidth}|}
\hline
\textbf{Сокращение} & \multicolumn{1}{c|}{\textbf{Расшифровка}} \\
\noalign{\hrule height 1.2pt}
\hline
API & Application Programming Interface (интерфейс прикладного программирования) \\
\hline
CI/CD & Continuous Integration / Continuous Delivery (непрерывная интеграция и непрерывная поставка) \\
\hline
DBMS / СУБД & Database Management System (система управления базами данных) \\
\hline
DevOps & Development and Operations (разработка и эксплуатация) \\
\hline
ERP & Enterprise Resource Planning (планирование ресурсов предприятия) \\
\hline
FMEA & Failure Mode and Effects Analysis (анализ видов и последствий отказов) \\
\hline
IDE & Integrated Development Environment (интегрированная среда разработки) \\
\hline
IS / ИС & Information System (информационная система) \\
\hline
ISO & International Organization for Standardization (Международная организация по стандартизации) \\
\hline
ITIL & Information Technology Infrastructure Library (библиотека инфраструктуры информационных технологий) \\
\hline
LC / ЖЦ & Life Cycle (жизненный цикл) \\
\hline
MTBF & Mean Time Between Failures (средняя наработка на отказ) \\
\hline
MTTR & Mean Time To Repair (среднее время восстановления) \\
\hline
OS & Operating System (операционная система) \\
\hline
QA & Quality Assurance (обеспечение качества) \\
\hline
RACI & Responsible, Accountable, Consulted, Informed (ответственный, подотчётный, консультируемый, информируемый) \\
\hline
RFP & Request for Proposal (запрос предложений) \\
\hline
ROI & Return on Investment (возврат на инвестиции) \\
\hline
RPO & Recovery Point Objective (целевая точка восстановления) \\
\hline
RTO & Recovery Time Objective (целевое время восстановления) \\
\hline
SLA & Service Level Agreement (соглашение об уровне обслуживания) \\
\hline
SLO & Service Level Objective (целевой показатель уровня обслуживания) \\
\hline
SOA & Service-Oriented Architecture (сервисно-ориентированная архитектура) \\
\hline
SRS & Software Requirements Specification (спецификация требований к программному обеспечению) \\
\hline
TCO & Total Cost of Ownership (совокупная стоимость владения) \\
\hline
UAT & User Acceptance Testing (приёмочное тестирование пользователями) \\
\hline
UI/UX & User Interface / User Experience (пользовательский интерфейс / пользовательский опыт) \\
\hline
VPN & Virtual Private Network (виртуальная частная сеть) \\
\hline
\end{tabular}
}

\clearpage
\nocite{*}
\printbibliography[heading=bibliography]

\end{document}
